% ============================= Discussion =============================
\section{Discussion and Limitations}
\label{sec:discussion}

\textbf{What is actually new?} That orientation is hard was already known from
VSR's original paper and re-confirmed by SITE. What we add is a
\emph{measured decomposition} of \emph{where} the failure lives, under one
protocol: (i) frozen-feature probes (linear, nonlinear, patch-vote,
object-grounded) quantify that the visual representation itself encodes
object-intrinsic direction only weakly; (ii) vision-side adaptation and explicit
two-stage decomposition are tested causally and fail (both significantly worse);
(iii) hard-negative training produces a coherence gain of $+40.8$ consistency
points at \emph{zero} accuracy change --- establishing logical consistency as a
separable, measurable failure dimension; (iv) a preregistered external validation
narrows the claim to object/direction orientation specifically. None of these
four facts existed before for this task.

\textbf{Only two model families.} SmolVLM2-2.2B and Qwen2-VL-7B were chosen for
depth, not breadth: every condition shares prompt, parser, splits, and metric
code, making the 15+ conditions mutually comparable. The SITE validation
provides an independent check on the 7B family; the 2B family shares the same
bottleneck profile, and per-relation orientation numbers for both are in
Table~\ref{tab:vsr}. We do not claim universality across architectures.

\textbf{Small orientation set ($n{=}137$).} This is a real limitation, addressed
three ways: Wilson 95\% CIs throughout; paired exact-McNemar tests (which do not
require large $n$); and an external, preregistered validation with $n{=}2{,}024$
orientation-vocabulary examples, where the same weakness appears at scale. We
report every per-relation cell with its $n$ so readers can discount small cells
(notably perpendicular, $n{=}12$).

\textbf{Is the SITE orientation subset cherry-picked?} The subset is
keyword-derived and \emph{non-official} --- we say so in Table~\ref{tab:site}.
It was frozen \emph{before} evaluation (example IDs committed with config hash
\texttt{28f4cc09887477af}), the keyword rule is fully specified in
App.~E, the source/category distribution is in
App.~E, and the preregistered headline subset (the official
spatial-relationship category) shows the contrast that makes the finding
meaningful: the same model, same prompts, same images --- strong on relations,
weak on orientation.

\textbf{Do the probes really establish a representation failure?} The probes
measure decodability by \emph{simple} readouts; a strong nonlinear decoder could
in principle do better (we tested MLPs --- no gain). They also share their
grounding boxes with the generative model, and region pooling is mean-based. We
therefore phrase the claim precisely: orientation is \emph{not robustly
decodable} by the readouts we tested, and the generative model's own residual
signal exceeds all of them, indicating the useful component lives in
multimodal/language interaction. The claim that matters for future work ---
``adapting the frozen encoder did not help either'' --- is a causal result, not a
probe result.

\textbf{Why should 66$\rightarrow$77 consistency matter if accuracy is flat?}
Because it is the one intervention that \emph{reliably changes what the model
believes}, and it shows accuracy is insensitive to a major qualitative failure:
the zero-shot model contradicts itself on 63\% of facing pairs --- below chance
--- while scoring 64\% accuracy. Reviewers of spatial reasoning have asked what
``understanding'' means for a model that answers a statement and its negation
both ``True''. Consistency is the measurable proxy we offer, and the
accuracy-consistency dissociation makes it possible to evaluate coherence
interventions independently.

\textbf{Many comparisons without correction.} We report all 23 paired tests
(complete battery in App.~D), exact binomial McNemar $p$-values,
discordant-pair counts, and per-cell $n$. No test was run and withheld; the
headline claims are each supported by multiple independent tests and by the
external validation. We do not apply a multiple-comparison correction because
the battery is small, fully enumerated, and the pattern (not any single
$p$-value) is the claim --- but we flag that readers should treat individual
$p$-values near 0.01 with the usual caution.

\textbf{Why no video?} The preregistered SITE protocol defers the video subset
(3{,}619 examples) and the VSR-trained LoRA on SITE until the zero-shot image
results were reviewed. Under decision-rule branch 2, the LoRA condition was not
justified: it does not fix orientation on VSR, and the claim was narrowed to
orientation-specific weakness. Video orientation (movement/navigation) is a
genuinely different capability (dynamic, egocentric direction) and an explicit
extension, not a gap in the current claim.

\textbf{Limitations.} Single seed per condition (no seed variance reported);
VSR subject/object parsing is heuristic; annotation of the 48 persistent failures
was performed by one annotator with a documented taxonomy; the 2B structured
prompt result is a single run; parallel/perp as strict complements is a soft
complement in reality (we handle this explicitly in the consistency analysis).
SITE media for the video subset requires the HF token to re-download; all other
artifacts are committed.

\textbf{Broader impact.} Understanding exactly which spatial capability resists
training --- object-intrinsic direction --- is directly relevant to robotics,
navigation, and document/figure understanding, where ``which way is it facing''
matters. Our decomposition template transfers to any VLM spatial claim; we also
release the complete artifact index and audit script so every number here can be
recomputed in minutes (\texttt{scripts/make\_paper\_figures.py}).
